% Source: https://artofproblemsolving.com/community/c3h2129717_twin_primes
% Stable identity: aops-topic-2129717-q01
The claim that $P_n-1$ and $P_n+1$ are always both prime is false. For example:
\begin{itemize}
\item $n=1$: $P_1-1=1$, which is not prime;
\item $n=4$: $P_4=2\cdot3\cdot5\cdot7=210$, and $P_4-1=209=11\cdot19$;
\item $n=6$: $P_6=30030$, and $P_6+1=30031=59\cdot509$.
\end{itemize}
What is true is that neither $P_n-1$ nor $P_n+1$ is divisible by any of $p_1,\dots,p_n$, since each leaves remainder $\pm1$ on division by every $p_i$. So every prime factor of $P_n\pm1$ is a prime \emph{other than} $p_1,\dots,p_n$. This is Euclid's argument that there are infinitely many primes, but it does not show that $P_n\pm1$ themselves are prime, so the argument proves nothing about twin primes. (Whether there are infinitely many twin primes is a famous open problem, the twin prime conjecture.)

So the error is the unjustified, and false, step $\boxed{P_n\pm1\text{ are both prime}}$.
